% ============================================================
% Section 3: Experimental Setup
% ============================================================
\section{Experimental Setup}
\label{sec:setup}

All experiments follow a standardized protocol designed for reproducibility.

\subsection{PUF Simulation}

We simulate XOR Arbiter PUFs using the standard linear delay model.
Each $k$-XOR APUF consists of $k$ independent 64-bit Arbiter PUFs whose
outputs are XORed.  Arbiter weights
$\mathbf{w}_i \in \mathbb{R}^{65}$ (64 delay differences plus one bias term)
are sampled i.i.d.\ from $\mathcal{N}(0, \alpha^2)$ with $\alpha = 0.05$.
All experiments use noise-free simulation ($\sigma = 0$) to isolate the
data scaling effect from hardware noise.

Interpose PUFs are simulated by inserting the upper-layer response into
a designated position (bit~0) of the lower-layer challenge, following the
specification of~\cite{wisiol2020splitting}.

\subsection{Attack Models}

We evaluate three families of models:

\paragraph{Standard MLP (baseline).} A 4-layer multi-layer perceptron with
hidden dimensions $2^{k-1}$, $2^{k}$, $2^{k-1}$, 1, Tanh activations in the
hidden layers and Sigmoid on the output, as used in prior
work~\cite{ruhrmair2010modeling,santikellur2023}.

\paragraph{Split MLP (NNModel\_fc).} A single 4-layer MLP with two independent
probability heads, the final response computed as $p_1(1-p_2) + p_2(1-p_1)$.
Parameter count for $k$-XOR is
$2^{k/2+1}$ in the hidden layers---roughly $7.7\times$ fewer than the
standard MLP at $k=6$.

\paragraph{Architecture variants (for independence study).} We test six
additional architectures at 8-XOR, 2M CRPs: Deep MLP (6-layer), Wide MLP
($k{=}9$ width, 4-layer), Deep+Wide ($k{=}9$, 6-layer), ReLU+Kaiming
initialization, Residual connections (2 skip connections), and
NoSigmoid+BCEWithLogits loss.  The largest variant has $1.08$M parameters.

\subsection{Training Protocol}

\begin{itemize}
\item \textbf{Optimizer}: Adam~\cite{kingma2015adam} with learning rate
  $\eta = 0.001$, $\beta_1 = 0.9$, $\beta_2 = 0.999$, cosine-annealed
  schedule~\cite{loshchilov2017sgdr} over $T_{\max} = 500$ epochs, and
  binary cross-entropy (BCE) loss.
\item \textbf{Batch size}: 4{,}096 (16{,}384 on RTX 4090 campaigns);
  up to 500 training epochs.
\item \textbf{Early stopping}: training stops at epoch $e$ if the best
  validation accuracy through $e$ remains below 52\% at $e > 150$ (55\% at
  $e > 200$ in the architecture-independence campaign), eliminating
  approximately 70\% of wasted compute on clearly failing runs without
  affecting successful runs.
\item \textbf{Data split}: 80/2/18 with held-out validation (later
  campaigns: 80/20 without validation, fixed 150 epochs); data are shuffled
  once before splitting.
\item \textbf{Data generation}: challenges are drawn uniformly from
  $\{-1, +1\}^{64}$; CRP generation is GPU-accelerated via vectorized
  batch computation.
\item \textbf{Hardware}: Experiments through 8-XOR ran on a single NVIDIA
  GeForce RTX 3080 Ti (12~GB VRAM) with CUDA 12.1 and PyTorch 2.1.2; the
  9-XOR measurements at and beyond $10^7$ CRPs and all 10-XOR measurements
  were run on an NVIDIA GeForce RTX 4090 (24~GB VRAM).
\item \textbf{Reproducibility}: each experiment generates a fresh PUF
  instance and CRP dataset; we do not fix random seeds across independent
  runs, capturing the full distribution of PUF hardness.
\end{itemize}

\subsection{Evaluation Metrics}

\begin{itemize}
\item \textbf{Test accuracy}: fraction of correctly predicted responses on
  the held-out test set (18\% of total CRPs), with model output rounded to
  $\{0, 1\}$.
\item \textbf{Success criterion}: test accuracy exceeds $90\%$.
\item \textbf{Success probability}: fraction of successful runs out of $n$
  independent trials, with $n$ between 3 and 30 depending on the
  configuration.  Confidence intervals are Wilson score intervals, except
  boundary cases ($p = 0$ or $p = 1$) and small-sample rows ($n \leq 5$),
  which use Clopper-Pearson exact intervals.
\item \textbf{$N_{50}$}: The CRP budget at which $P(\text{success}) = 50\%$,
  defined as the midpoint of a logistic regression of success probability
  on $\log_{10}(N)$, fitted over interior data points ($0 < p < 1$), with
  $R^2$ computed on the probability scale; linear interpolation between
  adjacent data points serves as a sensitivity check.  The 95\% CI on
  $N_{50}$ is obtained from the fit covariance.
\end{itemize}

Each $P(\text{success} \mid N, k)$ estimate mixes two sources of randomness:
PUF instance difficulty (each run draws a fresh instance) and SGD training
stochasticity.  The reported confidence intervals account for the combined
variation without attributing it to either source; a stratified separation
of the two would require a substantially larger experimental design and is
left to future work.

\subsection{Baseline Methods}

\paragraph{Logistic regression (LR).} We use scikit-learn's
\texttt{SGDClassifier} with log-loss, trained for 1{,}000 iterations with
early stopping tolerance $10^{-3}$.  LR is tested on all $k$-XOR
configurations where the data fits in CPU memory ($\leq 2$M CRPs).
For 8-XOR at $\geq 5$M CRPs, the feature matrix exceeds available RAM, and
LR is reported as inapplicable.
